\section{Setup}
\label{sec:setup}

\subsection{Corpus and categories}

LoCoMo~\cite{locomo} is ten conversations, each spanning multiple sessions over
a simulated months-long relationship, 369--689 turns apiece, with questions
annotated by category and by the turns that constitute their evidence. We use
two categories:

\begin{description}[leftmargin=1.2cm,style=nextline]
\item[Category 4, single-hop] \rQsingle{} questions answerable from one turn.
\item[Category 1, multi-hop] \rQmulti{} questions requiring evidence from several
  turns, frequently in different sessions.
\end{description}

Only questions carrying annotated evidence are scored. Retrieval is within a
conversation --- the correct scope, since a question about one relationship is
never answered by another's transcript.

\subsection{Retriever}

Deliberately plain, because the paper is about the retrieval \emph{policy} and a
sophisticated retriever would confound it: turns embedded with
\texttt{zenlm/zen-embedding-0.6b}, cosine similarity, brute force over every turn
in the conversation, no approximate index, no reranking, no query rewriting.
Brute force costs \rSearchMs{}\,ms per search over a whole conversation \meas{},
which is small enough that an index would buy nothing at this scale and would add
a recall variable we do not want.

Vectors are computed once and cached, so every policy in \S\ref{sec:policies}
scores against a byte-identical index. This matters: it means a difference
between policies cannot be an embedding difference.

\subsection{Two metrics, and why both are needed}

For a question with gold evidence set $E$ and retrieved set $R_k$ of size $k$:

\begin{align*}
\text{recall@}k\ \textbf{all} &= \Pr\left[E \subseteq R_k\right] \\
\text{recall@}k\ \textbf{any} &= \Pr\left[E \cap R_k \neq \emptyset\right]
\end{align*}

\textbf{all} is the one that matters for answering: a question needing three
facts is not answered by two of them, and a model handed two of three will
typically produce a confident wrong answer rather than an admission. \textbf{any}
is the diagnostic: it separates ``the retriever does not understand the
question'' from ``the retriever understands the question and cannot reach all of
the answer''. Reporting only one of the two is how the failure in this paper
stays hidden.
